\documentclass[letterpaper,10.5pt]{article}
\usepackage[empty]{fullpage}
\usepackage{titlesec}
\usepackage{enumitem}
\usepackage{hyperref}
\usepackage{xcolor}
\usepackage{tabularx}

\hypersetup{colorlinks=false,pdfborder={0 0 1},pdfborderstyle={/S/U/W 1},linkbordercolor=cyan,urlbordercolor=cyan}
\setlength{\parindent}{0pt}
\setlist[itemize]{itemsep=-3pt, topsep=0pt, leftmargin=15pt}
\titleformat{\section}{\large\bfseries}{}{0em}{}[\vspace{-1pt}\titlerule\vspace{2pt}]

\begin{document}

% --- FIXED HEADER TABLE ---
\begin{tabularx}{\textwidth}{@{}X r@{}}
{\LARGE \textbf{Yohei Ohata} \hspace{10pt} Curriculum Vitae} & \\ % Added missing '&' here
\noalign{\vspace{6pt}}
Center for Information and Neural Networks & \href{mailto:yohei.ohata.ee@nict.go.jp}{yohei.ohata.ee@nict.go.jp} \\
National Institute of Information and Communications Technology & \href{https://yohei.ohata.github.io}{yohei.ohata.github.io} \\
1-4, Yamadaoka, Suita, Osaka, Japan & \\
\end{tabularx}

\vspace{10pt}

\section*{Education}
\textbf{The University of Tokyo} \hfill April 2022 -- March 2024 \\
M.S. in Human Factors Engineering \\
Department of Human and Engineered Environmental Studies \\
Cumulative GPA: 4.0/4.0 \\
{\small \textit{Thesis}: Identifying Different Thought Processes for Each Learner Using Eye-Tracking and Hidden Markov Models} \\
{\small 1st in Class (1/48), Dean’s Award, Best Department Thesis Award} \\
\par\vspace{6pt}
\textbf{Tokyo University of Science} \hfill April 2018 -- March 2022 \\
B.E. in Electrical and Electronic Engineering \\
Department of Applied Electronics \\
Cumulative GPA: 3.5/4.0 \\
{\small \textit{Thesis}: A Method of Designing a Hilbert Transformer with Transmission Zeros at Specified Positions} \\

\par\vspace{6pt}
\section*{Employment}
\textbf{National Institute of Information and Communications Technology} \hfill January 2026 -- Present \\
Research Assistant \\
Center for Information and Neural Networks (CiNet) \\
Computational Social Neuroscience Group \\
\par\vspace{3pt}
\textbf{Hitachi, Ltd.,} \hfill April 2024 -- December 2025 \\
Research Engineer \\
R\&D Group \\
 \\
\par\vspace{3pt}
\section*{Publications}
\subsection*{Articles in Peer-Reviewed Journals}
\begin{enumerate}[label=\textbf{J\arabic*}, leftmargin=35pt]
  \item \textbf{Ohata, Y.}, Hachisuka, S., Kurita, K., \& Warisawa, S. (submitted). Identifying Different Thought Processes for Each Learner Using Eye-Tracking and Hidden Markov Models. \textit{}.
\end{enumerate}\vspace{2pt}
\section*{Research Experience}
\textbf{Research Assistant} \hfill January 2026 -- Present \\ 
National Institute of Information and Communications Technology \\ 
{\small {Laboratory: Computational Neuroscience \& Decision Making Laboratory}} \\ 
{\small \textit{Supervisor: Haruno Masahiko}} \\ 
\begin{itemize}
  \item Trained a toy Transformer model on random walks over a community-structured graph and visualized the internal representations as a Hopfield energy landscape.
  \item Elucidating shared and distinct computational principles underlying hippocampal memory and the attention mechanism.
\end{itemize}
\vspace{8pt}
\textbf{Research Assistant} \hfill September 2025 -- Present \\ 
The University of Tokyo \\ 
{\small {Laboratory: Amano \& Nakayama Lab}} \\ 
{\small \textit{Supervisor: Yamashita Ayumu, Amano Kaoru}} \\ 
\begin{itemize}
  \item Performed cross-subject temporal synchronizing on the fMRI dataset (AX-CPT, Cued Task Switching, Sternberg, Stroop) using BrainSync (MATLAB).
  \item Applied tensor decomposition to the DMCC55B dataset to decompose data into spatial, temporal, and subject-load modes (MATLAB).
  \item Quantified cognitive control performance by indexing RT and Accuracy via the Cognitive Control Index, extracting contributory components using conditional randomization testing (R).
\end{itemize}
\vspace{8pt}
\textbf{Research Assistant} \hfill April 2022 -- March 2024 \\ 
The University of Tokyo \\ 
{\small {Laboratory: Inovative Learning Creation Lab}} \\ 
{\small \textit{Supervisor: Hachisuka Satori, Warisawa Shin'ichi}} \\ 
\begin{itemize}
  \item Designed experimental protocols for both adult and elementary school participants, including the development of a web-based UI and logging system using HTML/CSS/JS.
  \item Managed high-fidelity eye-tracking data collection and processing using Tobii Pro 3 and Tobii Pro Lab.
  \item Performed advanced statistical analysis using Linear Mixed-Effects Models (LMM) in R.
  \item Modeled the temporal dynamics of scan paths using Hidden Markov Models (HMM) implemented in Python.
  \item Conceptualized research goals, authored manuscripts.
\end{itemize}
\vspace{8pt}
\section*{Professional Experience}
\section*{Teaching Experience}
\textbf{Sensing and Signal Processing} \hfill April 2023 -- June 2023 \\
Teaching Assistant, The University of Tokyo \\
\begin{itemize}
  \item Assisted in building signal processing systems using LabVIEW
\end{itemize}\vspace{5pt}
\section*{Honors and Awards}
IEEE CAS JJC Best Student Award \hfill 2021 \\ 
\section*{Languages}
\begin{itemize}[leftmargin=15pt]
  \item High Level in English
  \item Native in Japanese
\end{itemize}\vspace{5pt}
\section*{References}
\textbf{Haruno Masahiko, Ph.D.,} \\ 
Professor, National Institute of Information and Communications Technology \\ 
Email: \href{mailto:mharuno@nict.go.jp}{mharuno@nict.go.jp}\vspace{5pt} \\ 
\textbf{Ayumu Yamashita, Ph.D.,} \\ 
Associate Professor, Kobe University \\ 
Email: \href{mailto:ayumu722@people.kobe-u.ac.jp}{ayumu722@people.kobe-u.ac.jp}\vspace{5pt} \\ 
\textbf{Robert Mok, Ph.D.,} \\ 
Principal Investigator, National Institute of Information and Communications Technology \\ 
Email: \href{mailto:robmmok@gmail.com}{robmmok@gmail.com}\vspace{5pt} \\ 
\textbf{Satori Hachisuka, Ph.D.,} \\ 
Associate Professor, The University of Tokyo \\ 
Email: \href{mailto:hachisuka@edu.k.u-tokyo}{hachisuka@edu.k.u-tokyo}\vspace{5pt} \\ 
\end{document}